\chapter{效率技巧：像老手一样操作}

这一章收集的是“知道与不知道差别很大”的操作习惯。每条都不长，但能让你的建模速度翻倍。

\section{手势的效率用法}

\begin{center}
\small
\begin{tabularx}{\linewidth}{@{}lX@{}}
\toprule
\textbf{场景} & \textbf{技巧} \\
\midrule
手指挡住目标 & 长按进入\textbf{精密指针}：准星出现在手指上方 60 dp，抬起手指即完成拾取 \\
选不准小元素 & 先用两指放大，再轻触；拾取有优先级（顶点 12 dp → 边 10 dp → 面） \\
需要旋转草图视角 & 草图中单指用于绘制，\textbf{三指}才能旋转；或者点导航立方体 \\
误触清空选择 & 双击是“清除选择”，做精确选择时少用双击 \\
想右键 & \textbf{双指轻点}视图，弹出与桌面右键等价的菜单 \\
批量选对象 & \textbf{双指长按}进入框选，拖动矩形；配合选择过滤器只框面或只框对象 \\
退出当前工具 & 右下角的工具条上有 ✕（等同 Esc）；或按返回键 \\
\bottomrule
\end{tabularx}
\end{center}

\section{返回键 = 万能 Esc}

记住这条就够：\textbf{返回键按层关闭}，从最内层开始——数字键盘、文本输入、选择器、菜单、工具面板、Draft、
装配配合、任务面板、特征对话框、草图。全部关闭后应用退到后台，不会退出。所以“卡在某个模式里”时，
连按返回键总能回到正常状态。

\section{选择与选择过滤器}

\begin{itemize}
  \item 轻触是\textbf{累加选择}，再点同一元素取消；选择条上的 ✕ 移除单项；双击清空。
  \item 需要只选面或只选边：\menupath{工具} 里的选择过滤（或命令 \cmd{Part\_FaceSelection} 等），
        视图右下角会出现“仅选择面”的提示条，点 ✗ 清除过滤。
  \item 选择历史：\menupath{编辑 → 选择后退/前进}（\cmd{Std\_SelBack}/\cmd{Std\_SelForward}）可以在选择集之间来回。
\end{itemize}

\section{用表达式代替数字}

任何数值输入都能写表达式：数字键盘最上面直接输入，或以 \file{=} 开头。例如：

\begin{lstlisting}
=60 mm / 2          // 30 mm
=Sheet.Thickness*2  // 引用电子表格里的别名（第 11 章）
=10 mm + 5 mm       // 带单位的运算
\end{lstlisting}

好处是：改一个变量，所有引用它的尺寸一起变。复杂模型里，把关键尺寸集中放到电子表格（变量选项卡），
再用表达式引用，是最省心的做法。

\section{用 Tip 回放建模历史}

实体（Body）的 \ui{Tip} 标记决定“算到哪一步”。在树里长按某个特征，选 \ui{将 Tip 滚动到此处}，模型会回退到
那一步的状态：可以检查中间结果、在早期特征上做修改。再滚回最后的特征即可恢复完整模型。

配合 \ui{压缩}（Suppress）可以临时排除某个特征，对比“有它/没它”的结果；排除带来的错误也会一并消失。

\section{工具搜索与最近工具}

\begin{itemize}
  \item 工具面板顶部的 \ui{搜索所有工具} 是全局搜索：输入中文命令名（凸台）、英文名（pad）或命令 ID
        （PartDesign\_Pad）都能命中；
  \item 工具条右侧的“最近工具”保留 6 个，反复用的命令会留在手边；
  \item 长按工具面板里的按钮不会弹出说明，命令一词即是 FreeCAD 的命令名，可在附录 A 与网上文档对照。
\end{itemize}

\section{把常用面板放在“变量”选项卡}

电子表格不只是算数：把设计参数（长、宽、高、孔径、阵列数）写成一张表，建模时全部引用它。改一处，
整机更新。这是参数化设计的“驾驶舱”。

\section{性能与流畅度}

\begin{itemize}
  \item \textbf{隐藏不看的对象}：大装配里把无关构件隐藏（树里长按 → 隐藏），视图重绘会明显变快；
  \item \textbf{减少同时打开的对话框}：特征对话框实时预览，复杂模型上会持续计算；
  \item \textbf{分部件建模}：几百个特征的单一实体不如拆成几个实体再装配；
  \item \textbf{有限元网格}：移动端建议先用较粗网格试算，确认边界条件后再加密；求解时插上电源并让屏幕常亮；
  \item \textbf{存储}：运行时约 1 GB 常驻，剩余空间不足 2 GB 时系统可能杀掉后台应用。
\end{itemize}

\section{硬件键盘与鼠标}

平板接上键盘/鼠标后，输入框可直接键入，布局不重建（插入键盘不会打断当前操作）。鼠标左键=轻触，
右键=上下文菜单，滚轮=缩放（等价于双指捏合）。触控板的多指手势由系统处理。

\section{自动保存与恢复}

MobileCAD 每 60 秒（以及每次切到后台时）把当前文档写入恢复副本。若应用异常退出，下次启动会自动打开
恢复副本并提示“已恢复未保存的工作 from HH:mm”；这时请立刻“保存”到一个正式文件。正常保存后恢复副本会被删除。
